\chapter*{Abstract}
\addcontentsline{toc}{chapter}{Abstract}

This Technical Manual documents the Robot Explota Globos Web Platform, a system built to support participant registration, communication management, and operational control of the Robot Explota Globos tournament. The project is framed within Universidad Tecnologica de Pereira and the Semillero de Investigacion Kilby. The project responsibility is associated with Yurley Tatiana Tovar Martinez, and the development work was carried out by Moises Sanchez Naranjo and Tomas Bermudez Londono.

The platform is implemented as a modular Django project. The validated technical audit identified five local applications: \archivo{apps.core}, \archivo{apps.participants}, \archivo{apps.tournament}, \archivo{apps.invitations}, and \archivo{apps.common}. This separation reflects different responsibilities: public website, registration and attendance, competition logic, institutional communications, and reusable base models. The structure keeps public registration flows separate from internal tournament and communication operations, improving maintainability and reducing the risk of cross-domain changes.

The competition component is the most sensitive technical core. The system supports school and university registrations, attendance confirmation through UUID tokens, synchronization of confirmed registrations into operational teams, and construction of competitions by division. The tournament logic covers group stages, elimination brackets, repechage stages, progressive phases, manual phases, result correction, and final podium registration. This capability is mainly centralized in \archivo{apps/tournament/services.py}, which makes the main business rules easier to locate but also creates regression risk if changes are made without sufficient tests.

The communications module manages DOCX templates, XLSX recipients, document rendering, PDF conversion through LibreOffice, and email delivery using either SMTP or console backends depending on configuration. The module includes dry-run and controlled test modes to reduce the risk of accidental real deliveries. The frontend is based on Django Templates, shared CSS, and specific JavaScript files: \archivo{static/js/app.js} for public interactions and \archivo{static/js/control.js} for the internal panel, including AJAX operations with CSRF protection.

The technical audit received a score of 88/100 and was approved with observations. These observations do not block the manual writing process, but they guide the following technical chapters: detailing the \variable{DivisionCompetition.configuration} schema, expanding internal contracts for critical services, converting testing gaps into a QA matrix, avoiding invented deployment data, and marking production-related pending information as \MarcadorProduccionPendiente. This first block defines the context, objectives, scope, system overview, and implemented architecture that will support the remaining technical chapters.
